\chapter{Introduction}
\label{ch:intro}

Large language models (LLMs) have shown a remarkable ability to learn classification tasks without being specifically trained for them, and hence without changing their internal weights. By placing a few labelled examples (\emph{demonstrations}) in the prompt, the model can be seen to answer the query prompt sufficiently well \citep{brown2020language, dong2024survey}. This phenomenon, known formally as \emph{in-context learning}, has been widely used by both users and professionals during training, testing and inference, largely due to its ease of use and setup. However, due to its fragile nature, it brings about challenges and varying results depending on which demonstrations are shown, their format and order, and their label balance \citep{zhao2021calibrate, lu2022fantastically, min2022rethinking}. Hence, a body of work arose to investigate the \emph{selection} of demonstrations: by similarity to the query \citep{liu2022what}, by diversity and coverage \citep{levy2023diverse,gupta2023coverage}, or with retrievers trained on the model's own feedback \citep{rubin2022learning,li2023unified,ye2023compositional}.


To the extent of the work that follows here, there are two things in practice that have been overlooked by most of the current work out there. 
The first is \textbf{distribution shift}: the queries a deployed model sees are rarely drawn from the same distribution as the demonstrations.
The second is \textbf{prior conflict}: a model brings its own beliefs about what a feature means, and what the demonstrations teach can disagree with them. Work by \citet{wei2023larger} found that following in-context evidence that contradicts semantic priors emerges with scale, and models of 7 to 8B parameters follow contradicting demonstrations only partially (not being able to override their priors and learn from what the demonstrations teach) \citep{krishnakumar2025semantic,kossen2024incontext}.


\newpage
\section{The question}

This thesis aims to ask a single question:

\begin{quote}
\textit{Can the choice of demonstrations help a small language model follow in-context evidence over its own priors, under controlled distribution shift?}
\end{quote}

The question is then broken into four research sub-questions, each answered on Qwen2.5-7B-Instruct \citep{qwen2024qwen25} with $k = 8$ demonstrations.

\begin{description}[leftmargin=2.4em, style=nextline]
  \item[\RQ1] Does in-context learning degrade under covariate shift, spurious-feature reversal and semantic-prior conflict?
  \item[\RQ2] Do different selection strategies help under different kinds of shift?
  \item[\RQ3] Does a model of this size rely on the features that actually determine the label, and where in the model does its prior about them win when it conflicts with the demonstrations?
  \item[\RQ4] Does a learned selector (\sata) beat the hand-designed strategies?
\end{description}


It is difficult to answer these questions on real data, as real data and the benchmarks built on them mix several kinds of shift and give no ground truth for which features matter (\autoref{sec:bg-shift}). Therefore, the work that follows builds a \textbf{synthetic tabular testbed} (\autoref{ch:method}) where every task has three features that determine its label, a spurious ``shortcut'' feature that agrees with the label 80--90\% of the time, and distractors. The conditions in focus here, covariate shift, reversal of the shortcut and semantic-prior conflict, can each be switched on alone. Semantic-prior conflict is created by \emph{naming}: the same task is shown with neutral names ($f_0, \ldots, f_9$), with everyday names whose meaning agrees with the rule (\emph{aligned}), or with the antonyms of those names (\emph{flipped}), so that the model's semantic prior points the wrong way. Unlike conventional concept shift, the rule that generates the labels never changes; only what the names lead the model to expect does (\autoref{sec:environments}).

\newpage
\section{Approach and the story of this thesis}

The experiments here form a chain of sorts, and each chapter of Part~II aims to start from the limitation of the one preceding it.

\begin{enumerate}[leftmargin=2em]
  \item \textbf{Shift, names and the prior} (\autoref{ch:shift}; \RQ1--\RQ3).
        We measure how each shift hurts and whether seven hand-designed strategies help.
        The names dominate: none of the strategies we tested overcame flipped names, and the model's ``bigger numbers mean 1'' habit explains much of what remains.
        \emph{Limitation:} hand-designed heuristics see only the data, not what the model makes of it (information about the model's internal bias via its priors).
  \item \textbf{Learning to select} (\autoref{ch:sata}; \RQ4).
        We train a set transformer, \sata (Shift Aware Task Adapter), to choose demonstrations that teach the rule rather than the known shortcut.
        On a stand-in learner it is the only method with no weak environment, but on Qwen-Instruct it shows no significant advantage over balanced random demonstrations: it evens out the three environments at the cost of lower performance in distribution and under covariate shift. In an exploratory comparison, \sata does help the base model.
        \emph{Limitation:} behaviour alone cannot tell us why none of the demonstration sets we tried moves the instruction-tuned model.

  \item \textbf{Inside the model} (\autoref{ch:mechanism}; \RQ3).
        Using the logit lens, we saw that the answer appears to form in the last few layers and that the split between aligned and flipped names appears there; linear probes read the demonstrated rule only weakly from the hidden states, and steering towards the demonstrations shifts every answer without making it more correct. These tools show where the names' effect appears, not the full mechanism by which the model decides.
        \emph{Significant limitation:} our scale: one 7B model, eight numeric demonstrations, synthetic tasks.
\end{enumerate}

\section{Contributions}

\begin{enumerate}[leftmargin=2em]
  \item \textbf{A controlled testbed for in-context learning under shift and prior conflict}, with each shift isolated, ground-truth feature roles, a naming factor matched on measured prior strength, and a statistical protocol (within-cell \AUROC, a hierarchical bootstrap over tasks, Holm-corrected contrast families frozen before the data) (\autoref{ch:method}).
  \item \textbf{A measured account of how the 7B model here trades demonstrations against priors}: names outweigh every data shift, and the model's numeric prior explains much of the cost of reversing the shortcut (\autoref{ch:shift}).
  
  \item \sata, our demonstration selector, trained on ground-truth targets rather than on the model's scores, with an oracle check that at first caught a target teaching the shortcut backwards, and a pre-registered evaluation on the 24 evaluation tasks, which showed that \sata was not significantly better than the strategies on the instruction-tuned model at our small scale, and that its robustness there came at a cost in predictive performance (\autoref{ch:sata}).

  \item \textbf{Evidence from inside the model} on where the names' effect appears and how little of the demonstrated rule is visible there, and a negative result for single-direction steering (\autoref{ch:mechanism}).
\end{enumerate}

\section{Summary of findings}

The main takeaway is that, in our setting, the prior carried by the feature names influences the model's predictions more than the choice of demonstrations does: flipping three names costs 0.20 \AUROC, while the hand-designed strategies typically differ from balanced random demonstrations by only 0.025.
The findings are stated in bold throughout, numbered by chapter, and always for the conditions we tested.
The main ones are these below.

\begin{itemize}[leftmargin=1.5em]
  \item \textbf{Flipping feature names hurts more than any shift of the data, and far more than the choice of demonstrations helps}: 0.20 \AUROC, against 0.12 for reversing the shortcut and 0.02 (not significant) for covariate shift (\autoref{ch:shift}).
  \item \textbf{Strategies that neutralise the shortcut help under its reversal more than under covariate shift} (interaction +0.11, Holm $p = 0.024$), but no strategy helped significantly under flipped names (\autoref{ch:shift}).
  \item \textbf{A trained selector that is the most robust on a stand-in learner shows no significant advantage on the model}: it evens out the environments at a cost in predictive performance, and none of the six pre-registered selector contrasts favours it (\autoref{ch:sata}).
  \item \textbf{The logit lens places the answer in the last six of 28 layers, where the names' effect appears}; the model's prior is decodable from layer 2, while the demonstrated rule hardly is (\autoref{ch:mechanism}).
\end{itemize}